
\documentclass[12 pt]{paper}
\usepackage[utf8]{inputenc}
\usepackage[american]{babel}
\usepackage{paralist}%in order to enumerate inside a paragraph
%\usepackage[hyphens]{url}
\usepackage[colorlinks=true,linkcolor=blue,anchorcolor=black,citecolor=blue,filecolor=blue,menucolor=black,runcolor=red,urlcolor=blue]{hyperref} %to get hyperlinks
\usepackage{caption}
\usepackage{subcaption}
\usepackage{csquotes}
\newbox{\bigpicturebox}%to create collage 

\usepackage[style=apa,backend=biber]{biblatex}%to cite using apa7

\usepackage{geometry}%to change page specs
\usepackage{enumitem,varwidth}%image stuff
\usepackage[svgnames]{xcolor}%image stuff, so pictures are colored
%\usepackage{tikz}%image stuff
%\usetikzlibrary{shapes.geometric} %image stuff
%\usepackage{lmodern}
\usepackage{graphicx}%to follow the template from school
\usepackage{booktabs} %to follow the template from school
 \geometry{
	a4paper,
	left=25mm,
	top=25mm,
	right=25mm,
	bottom=25mm,
}
\parindent24pt
%to follow the template from school


\graphicspath{ {./bilder/} } %where the pictures are
\addbibresource{MasterLibrary1.bib}
\title{%
Ditigization of early 1900 postcards (of Axvall)  }

\author{Sonia Lindblom \and Sicheng Yan}

\date{June 2021}

\begin{document}
\begin{titlepage}
	\begin{center}
		\vspace*{1cm}
		
		{\Large \textbf{Collection Development: Aerial postcards}}
		
		\vspace{0.5cm}
	
		
		\vspace{1.5cm}
		
		\textbf{Sonia Lindblom \quad \quad Sicheng Yan}
		
		\vfill
		
		
		
		\vspace{0.8cm}
		
		\includegraphics[width=0.4\textwidth]{1200px-Högskolan_i_Borås_Logo.png}\\
		\vspace{0.4cm}
		Master's programme: Digital Library and Information Services\\
		Swedish School of Library an Information Science \\
		University of Borås\\
		07 June 2021
	
		
	\end{center}
\end{titlepage}

\section{Introduction}
This report of a digitization project of older postcards featuring, among others aerial photos of Axvall. -- insert dizitization in  bigger context ---. The aim of this project is to make the images available online to those interested in historical aerial photographs, for instance 
\begin{inparaenum}[i)]
	\item for geomorphology purposes \autocite[cf.][]{gomez2015}
	\item post card historical purposes
	\item history of the village
	\item military history
\end{inparaenum}
\section{Background}
Axvall is a village in Valle region, Skara Municipality, Väster Götaland county in Sweden. It was formerly a municipality itelsef, with railway station, mayor's office and other functions as well a big sanatorim, Stora Ekbergs Sanatorium. One proeminent part of Axvall's history is the housing of at most three military regments (, broden2021). One of them, Skarborgs regemente trained in Axevalla hed beteween 1696 and 1913 \autocite{frykmer1992}. The other two are reported by \textcite{broden2021}, Väster Götalands regemente I6 up to 1927 and Livhusararna K3 up to 1905. It is hard to determine what contributed to the end of the railway activity. However it is closely realted to the extinsion of the military trainings on Axevalla hed, the extinsion of the beer brewery and so on. Axevalla hed was of military interest though up to the middle 1980's.

\section{Methodological work}%this has to have about 1300 words)
\subsection{Selection of the material}

The postcards were in shoeboxes without any systematization. Some of them were bound with rubberbands in bunts of 20. The conditions varied from untouched cards to others played with by kids who painted on with crayon. They were them placed in acidic free plastic boxes and grouped in categories such as hospital, housing, aerial photographs and so on, separated with laminated folders, following the recommendations of ..... 
 
The aerial photographs were of major interest since among them there were cards  which were forbiden to be distributed due to restrictions from the Swedish Army. Since aerial photographs taken before 1950's are free to publish \autocite[][kap. 3, 3§]{lantm2018}, we decided that digitizing those images would be a good contribution for the history of Axvall. Further investigation revealed though that a photo that has been restricted before by is not included in the exception in \textcite{lantm2018}. We have applied for a new assessment from Swedish Army, and the images were asessed in the end of May by a Swedish Army official. We got verbal authorisation and a promise of a written document.
\subsection{Metadata}

\subsection{Image capture}
The process of image capture might have been one of the most bothersome since we wanted to make sure we followed the standards determined by the Swedish Kungliga biblioteket. The camera, though not so up to date, is a good hardware, a Cannon EOS 40D. The light temprerature was set to 5000 Kelvin; We sought for a neutral background in light gray an a static .... for each postcard. The images were captured with the help of the program EOS_Utility to make sure the camera settings and distance from the motifs were always the same. 

A color ruler was manufactured with the help of a sample color pallet from Perstorp which was then translated to LAB, conforing to .... standatds. 
\subsection{.... platform}% Or something like webdesign

\section {Critical Analysis}


\section{Conclusions}


\newpage

%\bibliographystyle{apacite}
%\bibliography{../../../MasterLibrary.bib}
\printbibliography


\end{document}


